\documentclass[10pt,a4paper]{amsart}
\usepackage[margin=19mm]{geometry}
\usepackage{amsmath,amssymb,mathtools}
\usepackage[scr=rsfs,cal=euler]{mathalfa}
\usepackage{xfrac}
\usepackage[unicode,hidelinks,bookmarks=false]{hyperref}
\newcommand{\ie}{i.e.\ }
\newcommand{\cf}{cf.\ }
\newcommand{\resp}{respectively\ }
\newcommand{\Subsection}{\subsection}
\providecommand{\nobreakdash}{\nobreak\hbox{-}}
\setlength{\emergencystretch}{2em}
\multlinegap0pt
\usepackage{mathtools}
\usepackage{amssymb}
\usepackage{subfig}
\usepackage{siunitx}
\usepackage{booktabs}
\usepackage{xcolor}
\usepackage{enumerate}
\newtheorem{lem}{Lemma}
\newtheorem{thm}{Theorem}
\newcommand{\CddB}{{C_{B2}}} % Dream const derivative
\newcommand{\CdvB}{{C_{B1}}} % Dream const vector
\newcommand{\HH}{\mathbb{H}}
\newcommand{\LL}{\mathcal{L}}
\newcommand{\LLinv}{{\mathcal{L}^{-1}}}
\DeclareMathOperator*{\argmin}{argmin}
\newcommand{\dG}{{\delta G}}
\newcommand{\etak}{{\eta_{k}}}
\newcommand{\norm}[2]{\left\| #1 \right\| _{#2}}
\newcommand{\pairing}[2]{\bigl\langle #1 , #2 \bigr\rangle}
\newcommand{\polN}{\mathbb N}
\newcommand{\s}{{\beta}}	% step size
\newcommand{\vk}{{v_{k}}}
\newcommand{\vkp}{{v_{k+1}}}
\title{A perturbed preconditioned gradient descent method for the unconstrained minimization of composite objectives}
\author{Jea-Hyun Park, Abner J. Salgado, Steven M. Wise}
\date{Source: arXiv:2512.19532v1 (CC BY 4.0)}
\begin{document}
\maketitle

\noindent\emph{Source and licence.} A perturbed preconditioned gradient descent method for the unconstrained minimization of composite objectives. Version arXiv:2512.19532v1 (CC BY 4.0), distributed by arXiv under the Creative Commons Attribution 4.0 International licence, http://creativecommons.org/licenses/by/4.0/, as recorded in the arXiv licence metadata for this exact version and shown on the abstract page https://arxiv.org/abs/2512.19532v1. Full licence notice: \texttt{licenses/arxiv-2512.19532-CC-BY-4.0.txt}.\par\smallskip

\noindent\textit{Editorial scope.} Counted unit: the article's thm thm:invariant-set (source0.tex lines 532--553). The article's complete original proof of it is delivered with it (source0.tex lines 555--606, one \texttt{proof} environment of 52 lines). No mathematics is written, repaired or re-derived: the statement and the proof are the article's own lines, in the article's own notation, and no retained line is substituted or re-flowed.  Retained uncounted as the source-local context the proof needs: lines 210--213; lines 382--384; lines 447--453; lines 449--451; lines 498--501; lines 503--505.  Nothing was omitted from any retained span, so the package carries no omission record.  No retained line contains a citation, so the wrapper carries no bibliography and no bibliography entry.  No retained line contains a figure environment, an image-inclusion command or drawing code, so no image is delivered and the package declares no asset dependency.  Editorial formatting only, in addition: an AMS \texttt{amsart} preamble that restates the article's own declarations and macros used by the retained lines, namely the environments \texttt{lem}, \texttt{thm} on separate counters, together with the standard packages those lines need.  The retained environments print in this document as Lemma 1, Theorem 1 in order of appearance, whereas the article prints the same items with the numbering of its own full document.  The complete native source of the article as distributed in the arXiv e-print for this exact version is bundled unchanged as \texttt{sources/191487/source0.tex}.
\begin{equation}
\label{eq:GisComposite}
  G(v) \coloneqq E(v) + F(v), \qquad \forall v \in \HH.
\end{equation}

	\begin{equation}\label{est:dual-lip}
		\frac{1}{\hat L_B}\norm{\delta G(v)-\delta G(u)}{\LLinv}^2 \le \pairing{\delta G(v)-\delta G(u)}{v-u}, \quad \forall u,v\in B.
	\end{equation}	

\begin{lem}[dual trap]\label{lem:dream}
	Let $G$ be $\mu$--strongly convex ($\mu>0$) locally Lipschitz smooth with respect to the $\LL$--norm. Then, for any $B\subset\HH$ bounded and convex, there are $\CdvB,\CddB>0$ such that, for any $u,v\in B$, we have
	\begin{equation}\label{est:dream}
		\pairing{\delta G(v)-\delta G(u)}{v-u} \ge \CdvB\norm{v-u}{\LL}^2+\CddB\norm{\delta G(v)-\delta G(u)}{\LLinv}^2,
	\end{equation}
	where $\CdvB \coloneqq(\mu^2 + 2\hat L_B)/(4\hat L_B+4\mu)$, $\CddB \coloneqq 1/(\mu+\hat L_B)$, and $\hat L_B$ is the constant from \eqref{est:dual-lip} applied to $v \mapsto G(v)-\frac \mu 4 \norm{v}{\LL}^2$.
\end{lem}

	\begin{equation}
		\pairing{\delta G(v)-\delta G(u)}{v-u} \ge \CdvB\norm{v-u}{\LL}^2+\CddB\norm{\delta G(v)-\delta G(u)}{\LLinv}^2,
	\end{equation}

	\begin{equation}
	\label{mthd:PPGD}
\vkp = \vk - \s\LLinv(\dG(\vk) + \etak) \quad \, k\ge 0,
	\end{equation}

	\begin{equation}\label{def:perturbation-PPGD}
\eta_k \coloneqq \eta_{v_k,\theta_k}\coloneqq   A(v_k;\theta_k) - \delta F(v_k),
	\end{equation}

	\begin{thm}[invariant set]
	\label{thm:invariant-set}
Let $G$ admit the decomposition \eqref{eq:GisComposite} and $\{v_k\}_{k\ge1} \subset \HH$ be obtained via the PPGD method \eqref{mthd:PPGD} with preconditioner $\LL$, initial guess $v_0\in \HH$, and with $\left\{A(v;\theta) \, \middle| \, \theta\in \Theta \right\}$ such that perturbations, defined in \eqref{def:perturbation-PPGD}, uniformly vanish. Set $u \coloneqq \argmin_{v\in\HH}G(v)$. Let $\CdvB>0$ and $\CddB>0$ be the constants appearing in \eqref{est:dream}. Define
	\[
d_0 \coloneqq \norm{v_0 - u}{\LL}, \quad \s_0 \coloneqq  \min\{\CddB, \CdvB^{-1} \} , \quad \epsilon_0 \coloneqq \left(\frac{1}{\CdvB} +\frac{5\CddB}{4} \right)^{-1} d_0^2 \CdvB.
	\]
If the step size is chosen so that
	\begin{equation}
	\label{est:sz-cond-invariant-set}
\s \in (0, \s_0],
	\end{equation}
and the solver parameters $\{\theta_k\}_{k\ge0}$ are chosen such that the perturbations $\eta_{k}$ satisfy, as guaranteed by the assumed uniform vanishing property,
	\begin{equation}
	\label{est:pert-cond-invariant-set}
\norm{\eta_{k}}{\LLinv}^2 \le \epsilon_0.
	\end{equation}
Then, for all $k \in \polN$, the iterates of PPGD satisfy
	\begin{align}
	\label{obj:invariant-set}
\vk \in B \coloneqq \left\{ v\in\HH \, \middle| \, \norm{v-u}{\LL} \le d_0  \right\}.
	\end{align}
	\end{thm}

	\begin{proof}
  Set, for $k \in \polN_0$, $d_k \coloneqq \norm{v_k - u}{\LL}$. We will show inductively that $d_k \leq d_0$, with the base for induction being trivial. Fix $k \in \polN$ and suppose $\vk\in B$. The assumptions on $G$ allow us to invoke Lemma~\ref{lem:dream} over $B$ so that
	\begin{equation}\label{misc10}
		\pairing{\dG(\vk)}{\vk-u} \ge \CdvB d_{k}^2 + \CddB \norm{\dG(\vk)}{\LLinv}^2.
	\end{equation}
	In addition, we trivially have
	\begin{equation}\label{misc09}
		\pairing{\etak}{\vk -u}
		\le	\frac{1}{2\CdvB} \norm{\etak}{\LLinv}^2+  \frac{\CdvB}{2}d_{k}^2
	\end{equation}
	and
	\begin{equation}\label{misc08}
		\pairing{\dG(\vk)}{\LLinv\etak}
		\le 
		\norm{\dG(\vk)}{\LLinv}^2+ \frac{1}{4} \norm{\etak}{\LLinv}^2.
	\end{equation}
	
	Subtract $u$ from \eqref{mthd:PPGD} and take $\LL$--norms. The choice of $\s$, \eqref{misc10}, \eqref{misc09}, and \eqref{misc08} yield,
	\begin{equation}
	\label{misc13}
	  \begin{aligned}
      d_{k+1}^2 &=\norm{\vk-u -\s\LLinv \delta G(\vk)-\s\LLinv\etak}{\LL}^2
      \\
      &= d_{k}^2 + \s^2\norm{\delta G(\vk)}{\LLinv}^2 + \s^2\norm{\etak}{\LLinv}^2 -2\s\pairing{\dG(\vk)}{\vk-u}
	\\
& \quad - 2\s\pairing{\etak}{\vk-u} +\s^2\pairing{\dG(\vk)}{\LLinv\etak}
      \\
      &\le (1-2\s \CdvB) d_{k}^2 - \s(2\CddB -\s)\norm{\delta G(\vk)}{\LLinv}^2 + \s^2\norm{\etak}{\LLinv}^2
	\\
& \quad - 2\s\pairing{\etak}{\vk-u} +\s^2\pairing{\dG(\vk)}{\LLinv\etak}
      \\
      &\le \left(1-\s\CdvB \right)d_{k}^2 - 2\s\left(\CddB -\s\right)\norm{\delta G(\vk)}{\LLinv}^2
	\\
& \quad +\s\left(\frac{1}{\CdvB} +\frac 5 4 \s\right)\norm{\etak}{\LLinv}^2
      \\
      &\le \left(1-\s\CdvB \right)d_{k}^2  + \s \left(\frac{1}{\CdvB} +\frac{5\CddB}{4} \right)\norm{\etak}{\LLinv}^2,
	  \end{aligned}
	\end{equation}
	where $\s \le \CddB$ was used in the last step and $0<\s \le \CdvB^{-1}$ ensures $0 \le \left(1-\s\CdvB \right) <1$.

	Now, since the perturbations uniformly vanish, for the bounded set $B$, we can choose $\theta_k$ such that \eqref{est:pert-cond-invariant-set} holds.
	Thus, we continue estimate \eqref{misc13} by using \eqref{est:pert-cond-invariant-set} and $\s \le \CdvB^{-1}$
	\[
		d_{k+1}^2 
		\le \left(1 - \s\CdvB \right) d_k^2 
		+
		d_0^2 \s\CdvB
		\le \left(1-\s\CdvB \right)d_0^2 +d_0^2 \s\CdvB
		= d_0^2.
	\]
	Therefore, we have $d_{k+1} \in B$, which completes the proof.
\end{proof}

\end{document}
